\section{Method}\label{sec:method}

Every rule in this section is taken from the pre-specification, which was frozen on
\FrozenDate{} before any code of the server existed. Its text was never edited after the freeze.
Readings of it made during engineering, values set by its rules, and the decisions of the author
and the study's coordinator went into its dated revision log. That log is part of the frozen file.
Where this section states such a reading, it says so.

\subsection{Families and Hypotheses}\label{sec:hypotheses}

The hypotheses form three families, and no contrast crosses families. No proxy is compared with
the server's in-process mode, and no result of the mechanism family supports a claim of cost.
Table~\ref{tab:hypotheses} lists the hypotheses.

\begin{table}[H]
\caption{The hypotheses, as frozen. A ratio's numerator is the first arm named. ``Holm'' marks a
family tested with Holm's step-down procedure at family-wise $\alpha = \FwAlpha$, one-sided.
\label{tab:hypotheses}}
\centering\small
\begin{tabularx}{\textwidth}{@{}l X l l@{}}
\toprule
\textbf{Id} & \textbf{Claim} & \textbf{Bound} & \textbf{Cells} \\
\midrule
C1 & Churn throughput (connections per second, WL1), one-port / dedicated, is equivalent & within
$[\MarginLow, \MarginHigh]$ & C1 to C3: \\
C2 & Keep-alive throughput (requests per second, WL3) is equivalent & within $[\MarginLow,
\MarginHigh]$ & at most \CostCellsMax, \\
C3 & Median time to first response byte at a fixed load (WL2) is equivalent & within
$[\MarginLow, \MarginHigh]$ & Holm \\
\midrule
B1 & The outcome and transcript of every hard case equal the frozen table & none deviates &
deterministic \\
B2 & Timed events are never early and are handled in the first pass after their deadline;
pending memory stays within its bounds & none deviates & deterministic \\
B3 & Counted footprint per pending connection, competitor / server, exceeds the margin & median
$Q > \FootprintMargin$ & \FamilyBCells, Holm \\
\midrule
M1 & The default detection mode is cheaper: churn throughput, default / other & median $>
\MechBound$ & \MechOneCells \\
M2 & In-process dispatch costs less CPU than relay: CPU per connection (WL6), in-process / relay &
median $< \MechBound$ & \MechTwoCells \\
M3 & The server's relay is faster than each proxy: churn throughput, server / proxy & median $>
\MechBound$ & \MechThreeCells \\
\bottomrule
\end{tabularx}
\end{table}

\runin{Cost family.} The server in one-port mode is compared with itself in dedicated mode: the
same binary, in-process dispatch, the default detection mode and one protocol per cell. The cells
cross C1, C2 and C3 with HTTP/1.1, h2c, TLS with HTTP/1.1 and MQTT. They run on epoll and
io\_uring on L and on IOCP on W, at most \CostCellsMax{} in all. The confirmatory family is the subset the A/A
pilot resolves (Section~\ref{sec:pilot}). Its size is \mC{} = \CostM. Stated before the run: the
TLS handshake dominates a TLS connection under churn. So the TLS cells have little power to show
the cost of a check of \TlsNeed{} bytes, and they are not read as evidence about it.

\runin{Robustness family.} B1 and B2 are deterministic and outside Holm. B1 runs every hard case of
the frozen table on every backend, in both detection modes, with every dispatch mode that
applies, \CaseReps{} times. The replicates exercise the timing of arrival. One deviation fails
B1. A connection classified other than as its script's protocol, in any measured window of any
family, also fails B1. B2 holds five bounds.
\begin{enumerate}[label=(\alph*)]
\item No fallback, close or deadline event happens before its deadline on the server's clock.
\item Each deadline is handled in the first loop pass whose wait returned at or after it.
\item Each connection is classified in the pass that received its deciding byte.
\item No pending connection holds more user-space payload than its mode's bound, and one that has
received no byte holds no data buffer.
\item A half-close before a decision closes the connection in the pass that observes it.
\end{enumerate}
The lateness of timed events and the decision
latency are reported per backend as distributions, and not tested.

B3 compares the counted footprint per pending connection, $W$ (Section~\ref{sec:footprint}), of
each competitor as configured with the server's. The author fixed its margin before any window
of B3 ran: a cell counts only if $W_{\mathrm{comp}}/W_{\mathrm{srv}}$ is at least
\FootprintMargin. The session statistic is $Q = W_{\mathrm{comp}}/W_{\mathrm{srv}}$, each arm's $W$
the mean of its two windows. A session whose $W_{\mathrm{srv}}$ is zero or less takes \Qzero,
which counts against. The server runs in relay mode against the proxies and in-process against
the libraries, in its default detection mode, with every detection timer at \PendTimerS{}~s. The
\FamilyBCells{} Holm cells cover two cases, each on epoll and io\_uring. The silent case runs
against nginx, HAProxy, Envoy, sslh-ev and hyper-util (\SilentHolmCells{} cells). The
partial-ClientHello case runs against the first four (\PartialHolmCells{} cells), since hyper-util
serves no TLS. Netty, Jetty,
caddy-l4 and cmux give \FamilyBDescriptive{} more cells, reported against the same bound and
deciding nothing: they measure a runtime's collector as much as a detector.

\runin{Mechanism family.} M1 compares the two detection modes for HTTP/1.1 and h2c churn on all
three backends, in the direction rule~E fixes (Section~\ref{sec:order}). On IOCP the peek mode
includes the switch to replay, so M1 there measures the limits of the API as much as the idea.
M2 compares the CPU time of the front and the back end together, per connection, for HTTP/1.1
and TLS on epoll and io\_uring. The relay arm passes TLS through to a back end that terminates
it, so both arms make one handshake per connection. M3 compares the server's relay with each of
the five proxies, each with one front core, in front of the stub. It covers TLS routed by its
server name and plaintext HTTP/1.1. The server runs on epoll, the accept model all five use. Each system holds
exactly two routes: TLS by server name to one stub port, and everything else to the HTTP/1.1 stub
port. What each examines to choose differs: nginx tells TLS from not TLS, while the server runs
its whole table.

\subsection{Workloads and Metrics}\label{sec:workloads}

\runin{WL1, churn in closed loop.} $C = \ConnSlots$ connection slots per server core. Each slot
repeats connect, one exchange and close (Table~\ref{tab:exchanges}). The metric is connections
per second: the exchanges completed without error per second of the measured window.

\begin{table}[H]
\caption{The exchanges of WL1.\label{tab:exchanges}}
\centering\small
\begin{tabularx}{\textwidth}{@{}l X l@{}}
\toprule
\textbf{Protocol} & \textbf{Exchange} & \textbf{Closes first} \\
\midrule
HTTP/1.1 & \texttt{GET} with \texttt{Connection: close}; read the response & server \\
h2c & preface, SETTINGS, one HEADERS frame with END\_STREAM; read to END\_STREAM; GOAWAY & client \\
TLS (HTTP/1.1) & full handshake, then the HTTP/1.1 exchange & server \\
MQTT & CONNECT; read CONNACK; DISCONNECT & client \\
SSH (secondary) & send the identification line; read the server's & server \\
TLS against the stub (M3) & one recorded ClientHello, the same bytes on every connection; read the
stub's \BodyBytes{} bytes to the end & server \\
\bottomrule
\end{tabularx}
\end{table}

\runin{WL2, churn in open loop.} Exchanges fall due at a rate $\lambda$, \RateFrac{} times a median
taken over the A/A pilot's sessions. Each session contributes its mean connections per second in
the C1 pilot cell with the same protocol and backend. The factor is a design choice, below
saturation. The time to first response byte ends at the first byte the server sends on the
connection, as the generator receives it. In open loop it runs from the exchange's due time, so a
stall of the generator cannot hide. The metric is its median. A window in which fewer than
\OpenDonePct\% of the exchanges due completed is invalid.

\runin{WL3, keep-alive.} $C$ connections are established before the window, each with one request
in flight. The request is an HTTP/1.1 \texttt{GET}, one h2 stream at a time, MQTT PINGREQ after one
CONNECT, or an HTTP/1.1 \texttt{GET} over one TLS connection. The metric is requests per second. A connection's
setup is not a request.

\runin{WL4 to WL6.} WL4 is secondary: the server's CPU time over the window per exchange, and its
resident memory at the end and at its peak. WL5 is secondary: the counters of
Section~\ref{sec:counters} per connection. WL6 is M2's metric: the CPU time of the front and the
back end together, per connection completed, at an open-loop rate. The rate of each M2 cell is
\RateFrac{} times the smaller of its two arms' median closed-loop connections per second. Those
medians come from \RateSessions{} development sessions on the frozen binary after the pilot. The relay
arm's back end runs in a \texttt{SO\_REUSEPORT} group of two workers, a design choice the revision
log records with its reason.

\runin{WL8, hard cases.} For each case, the outcome, the transcript and the decision time against
\opcase's write times. The table holds \HardCases{} cases. They cover whole, split and dripped
signatures for every protocol, and slow drips still incomplete at \Tdec. They cover silent clients
with and without a fallback, and a client whose first byte arrives just after or just before
\Tfb. PROXY headers come whole, split, missing, incomplete, too long and followed by each protocol.
The rest are bytes that rule out every matcher, half-closes and resets during detection, and a
Redis inline command, which is classified as HTTP/1.1 and answered with \texttt{400}.

\subsubsection{WL7, the Footprint per Pending Connection}\label{sec:footprint}

A connection is pending while the client's bytes are incomplete, whatever the system has decided
from them so far. Every B3 window lasts \PendWindowS{}~s from the first connect, in phases that
are design choices.
\begin{itemize}
\item Before \Tzero: fresh processes, the probe and a baseline sample.
\item Opening, from \Tzero{} to at most \PendOpenS{}~s: \opcase{} opens \NPend{} connections in
batches of \PendBatch, one batch every \PendBatchMs{}~ms. A batch is half of sslh's hard-coded
listen backlog of \SslhBacklog, the smallest among the systems; the pace opens \PendRate{}
connections per second. Each connection then sends nothing (the silent case), or the TLS record
header of the recorded ClientHello, of \HelloLen{} bytes, followed by the first half of it (the
partial-ClientHello case).
\item Settling, to $t = \PendSettleS$~s.
\item The first sample at $t = \SampleOneS$~s and the second at $t = \SampleTwoS$~s; the second
is the value used.
\item Closing, to $t = \PendWindowS$~s: every connection is closed by reset, so no socket enters
TIME-WAIT.
\end{itemize}
Each sample reads three parts, per pending connection. $U$ is the growth over the baseline of the
summed resident memory of the system's processes; the stub behind a relay holds no pending
connection and is left out. $K_q$ is the mean receive-queue allocation (the \texttt{r} field of
\texttt{ss -tm}) over the system's accepted sockets. It sums the true size of the buffers queued
for reading, which holds their payload and their headers. $K_s$ is the growth over the baseline
of the kernel's slab memory, less the growth of the socket-buffer caches. It is host-wide and
holds both loopback ends; the client ends are \opcase's, the same for every system.

A queued socket buffer's header is itself a slab object, so the slab total holds it while $K_q$
charges it too. In the Linux source at tag \KernelTag~\cite{linux2026v72}, those headers come from
three caches. The growth of the three is subtracted from the slab's growth, so each queued buffer
is counted once, in $K_q$. On L one of the three is merged into a shared cache, whose co-tenants
are \name{pool\_workqueue} and \SgPool, as read on L. The whole growth of the shared cache is
subtracted. Two residuals remain, and both sit with a system that leaves bytes queued. A header
too large for its cache is counted twice. Memory of the caches that no receive queue is charged
for is removed. The first raises that system's $W$ and the second lowers it. In the silent case
neither applies.

The holder \ophold{} measures \Kbase: it holds \NPend{} silent connections in the same layout, and
\Kbase{} is the median of its $K_s$ over \KBaseWindows{} windows, corrected the same way. A
system's footprint is $T = U + K_q + (K_s - \Kbase)$, and its counted footprint is $W = T + \Kbase
= U + K_q + K_s$. $W$ is the footprint of each system as configured, not of its detection design
alone, and it is not the whole system's. Outside it are the reserved but unused part of a socket's
receive memory, and kernel memory outside the slab that no receive queue is charged for. The second
residual and the co-tenants' growth of a merged cache are outside it too, and the first residual
can count bytes twice. \Kbase{} cancels from every statistic of B3. It is reported as the share of each footprint
that every system pays alike.

Every process raises its soft limit of open files to the hard limit. nginx's worker connections
and open files, and HAProxy's global connection limit, are twice \NPend; sslh's connection limit is
unset; Envoy and the harnesses set no downstream limit. Listen backlogs are \NPend{} wherever a
system has a setting, else the kernel's \texttt{net.core.somaxconn}. Runtimes with a collector run
it before each sample. The JVM systems get a garbage collection through \texttt{jcmd}, and the
cmux harness frees its heap to the system. caddy-l4 serves a heap profile, which collects first. Each reading,
the baseline included, follows the same step.

\subsection{Competitors and Their Configurations}\label{sec:competitors}

Table~\ref{tab:competitors} lists the competitors and their pins. Every configuration file is in
the paper's repository, every line with its reason. The rule is each project's documented best
configuration for the compared feature, with logging of individual connections off.

\begin{table}[H]
\caption{Competitors, pins and families.\label{tab:competitors}}
\centering\small
\begin{tabularx}{\textwidth}{@{}l X l@{}}
\toprule
\textbf{System} & \textbf{Pin} & \textbf{Families} \\
\midrule
nginx, stream module~\cite{nginx2026} & \NginxVersion, built from the release & B3, M3 \\
HAProxy~\cite{haproxy2026} & \HaproxyVersion & B3, M3 \\
Envoy~\cite{envoy2026} & \EnvoyVersion, the release binary & B3, M3 \\
caddy-l4~\cite{caddyl42026,caddy2026} & \CaddyLfVersion{} on Caddy \CaddyVersion, built with
xcaddy \XcaddyVersion{} and Go~\GoVersion & B3 (descriptive), M3 \\
sslh, the \texttt{sslh-ev} binary~\cite{sslh2026} & \SslhVersion & B3, M3 \\
Netty~\cite{netty2026} & \NettyVersion, on the Temurin JDK \JdkVersion & B3 (descriptive) \\
Jetty~\cite{jetty2026} & \JettyVersion, on the same JDK & B3 (descriptive) \\
cmux~\cite{cmux2021} & \CmuxVersion{} on Go~\GoVersion & B3 (descriptive) \\
hyper-util~\cite{hyperutil2026} & \HyperUtilVersion, with Rust \RustVersion & B3 \\
\bottomrule
\end{tabularx}
\end{table}

Each proxy has three configurations. The \emph{cases} configuration holds every route its features
cover, the matched timers and the fallback where it has one. The \emph{M3} configuration holds
exactly M3's two routes. The \emph{B3} configuration adds WL7's limits, the \PendTimerS{}~s timer
and the buffer settings below to the M3 configuration. Each
library harness has two, \emph{cases} and \emph{B3}. In M3 every front gets one core; in B3 every
competitor gets the server's core count. The detection timers of M3 are matched to the server's
\TimerS{}~s.

caddy-l4 is pinned by its release and by its commit, \hex{\CaddyLfCommit}.

\runin{Proxies.} nginx runs one worker on the front core. It uses the event method nginx chooses,
which on Linux is epoll with peer half-close reporting, and so its peek path. It routes with
\texttt{ssl\_preread} and accepts with \texttt{multi\_accept on}. HAProxy runs in TCP mode with
one thread, from the affinity set at start. It accepts a ClientHello by its hello type and plain
HTTP/1.1 at once, and routes by server name. Its client, server and connect timeouts stay unset,
its default. HAProxy gets \texttt{option splice-auto} only where the server's relay in the same
cell splices. Envoy follows its benchmarking advice: the release binary, one worker thread,
circuit breaking disabled and filter chains only for the compared features. Its default
statistics stay on, as the server's counters do. caddy-l4 runs with one Go processor, from the
affinity mask. sslh runs its \texttt{sslh-ev} binary with per-connection
logging off and numeric addresses.

\runin{Libraries.} The Netty harness follows Netty's port unification example with its built-in
TLS, h2c and PROXY handlers, on the native epoll transport. The Jetty harness uses its detecting
connection factory with TLS and PROXY, then HTTP/1.1 with the h2c upgrade. Both JVMs run the G1
collector, selected explicitly, with the JDK guide's default heap bounds written out, on IPv4
sockets. The cmux harness uses cmux's matchers for HTTP/2, HTTP/1, TLS and \texttt{SSH-}, with a
match-all last for the fallback. The hyper-util harness serves HTTP/1.1 and h2c on a tokio runtime
with one worker thread; it has no detection timer. Where a library terminates TLS, the settings of
Section~\ref{sec:handlers} apply as far as it allows.

\runin{Buffers in B3.} A default that sizes what a pending connection holds is changed in two cases
only. Either the project documents a value for deployments that face many untrusted connections,
or it documents that the buffer grows on demand, so a smaller start refuses nothing. So Envoy's
per-connection buffer limit is \EnvoyConnLimit{} bytes on the listener and the cluster, as its edge
guidance asks of TCP proxies. Its TLS inspector starts with \EnvoyInitialRead{} bytes, the
smallest value it accepts. HAProxy sets small buffers in its back ends, which a
pending connection does not use. Every other default is kept.

\subsection{Hosts and Procedures}\label{sec:hosts}

\runin{Hosts.} L is a Linux host with an \HostLCpu{} processor: \HostLCpus{} logical CPUs, with
\HostLSmt{} hardware threads per core. It runs Arch Linux with kernel \KernelL, clang
\ClangVersion{} and Go~\GoVersion. W is a Windows host with an \HostWCpu{} processor, \HostWCores{}
cores and \HostWCpus{} logical CPUs, running \HostWOs{} (build \HostWBuild) and MSVC
\MsvcVersion. Every run uses the loopback interface. The pre-specification cites Linux source lines
at tag \KernelTag; the revision log checked each cited line against the kernel the runs use and
found none changed in content.

\runin{Placement.} The hardware sibling of every used core stays idle. On L, CPUs \PlaceSysA{} and
\PlaceSysB{} run the system and the driver. In in-process cells the server runs on CPU
\PlaceInprocServer{} and \opgen{} on CPUs \PlaceInprocGenLo{} to \PlaceInprocGenHi. In hand-off
cells the front runs on CPU \PlaceFront, the back end on CPUs \PlaceBackA{} and \PlaceBackB, and
\opgen{} on CPUs \PlaceHandGenLo{} to \PlaceHandGenHi. In B3 and the hard cases the system or
\ophold{} runs on CPU \PlaceCaseSystem, with its back end on CPUs \PlaceCaseBackA{} and
\PlaceCaseBackB. \opcase{} then runs on CPUs \PlaceCaseLo{} to \PlaceCaseHi. The secondary cells with
two server cores run the server on CPUs \PlaceTwoA{} and \PlaceTwoB{} and \opgen{} on CPUs
\PlaceTwoGenLo{} to \PlaceTwoGenHi. On W the server runs on one logical CPU of the last physical
core, its sibling idle, and \opgen{} on the cores between core zero and the server's. The
placement follows the earlier study of the same author~\cite{tsvetanov2026routing}.

\runin{L.} Every run on L is a lab job under a lock. The job fixes the processor's state (boost
off, the performance governor, mains power) and holds each CPU's minimum frequency at its
maximum. It sets transparent huge pages to \texttt{madvise}, and it exempts loopback traffic from
connection tracking. The revision log records the reason for each. It
restores every setting when it ends, also on a signal. Each session's fingerprint records the
kernel, the settings before the job and during it.

\runin{W.} For each session the harness switches to a power plan of its own. It is a copy of
Windows' High performance plan with boost off, the processor state held at its maximum, and no
core parking. It switches back when the session ends, also on failure, without
administrator rights. The timer resolution stays at Windows' default and is recorded per window. A
quiet check before each job refuses a busy host. The frozen procedure also names a frequency
counter, tested against a known load. If no counter passes the test, W's windows are validated
without the frequency rule, and the paper says so (Section~\ref{sec:threats}).
\ForAlex{Section~9.1 of the frozen text asks for your approval of W's procedure. Record it, and
the counter's outcome, in the revision log at the code freeze.}

\subsection{Sessions and Validity}\label{sec:sessions}

A session is four windows of a cell's two arms in mirrored order, $XYYX$; which arm is $X$ is
drawn per session. Each window starts fresh processes, checks one exchange of each protocol the
cell uses (the probe), then runs a warm-up of \WarmupS{}~s and a measured window of
\MeasureS{}~s. The design follows the earlier study~\cite{tsvetanov2026routing}. B3 windows have
the layout of WL7. An arm's value in a session is the mean of its two windows. The session's
statistic is the ratio of the two arms' values, in the direction each hypothesis states. Sessions
are separate process lifetimes and are treated as independent. An analysis clustered by lab job
is reported beside every family and decides nothing.

A window is invalid, and is listed with its reason, if any of these holds:
\begin{itemize}
\item the server, the back end or the probe failed, no exchange completed, or any connect failed;
\item on L, the share of failed exchanges exceeded \ErrorSharePct\%, or the generator's CPUs were
more than \GenBusyPct\% busy in a closed-loop cell;
\item on L, the mean CPU frequency drifted more than \MhzDriftPct\% from the session's value;
\item on W, a rule of W's procedure that can be computed there failed;
\item any connection was classified other than as its script's protocol (also a failure of B1);
\item in a hand-off cell, a back-end core was more than \BackendBusyPct\% busy;
\item in open loop, fewer than \OpenDonePct\% of the exchanges due completed;
\item in the cost family, M1, M2 and B3, the kernel's listen overflow or drop counter changed
during the window;
\item in a B3 or \ophold{} window, $W$ moved between the two samples by more than its settling
bound;
\item in such a window, the host's count of TIME-WAIT sockets at either sample differed from the
baseline's, or fewer than \NPend{} of the system's accepted sockets were established;
\item in such a window, a socket-buffer cache could not be found or read.
\end{itemize}
The settling bound is the larger of \SettlePct\% of the second sample's $W$ and \SettleAbs{} bytes
per pending connection. In M3 the proxies keep their own backlogs, so their listen overflows are
reported beside each cell instead.
TIME-WAIT sockets are slab objects that last \TimeWaitS{}~s, so the first B3 or \ophold{} window of
a lab job starts at least that long after any other window ended.

A session with an invalid window is discarded and run again at the end of its family's order, at
most \RerunCap{} times per cell. A cell with fewer than $R$ valid sessions cannot pass and enters
Holm's procedure with \PIsOne{}. Beside every decision the paper reports the cell's invalid windows
per arm, with their reasons.

\subsection{Statistics}\label{sec:stats}

Every tested cell is decided by two computations, both fixed before the run. A cell passes only if
it passes under both; if they disagree, both are reported and the cell does not pass. The cell's
statistic is the median of its $R$ session values.
\begin{itemize}
\item \textbf{Clustered BCa interval.} \BootResamples{} resamples of the sessions with replacement,
the session as the cluster~\cite{efron1987bca,efron1993bootstrap}. The bias correction counts ties
as half, and the acceleration comes from the jackknife over sessions. The share of resamples beyond
a bound is clipped to \ClipRange, with $B = \BootResamples$. The one-sided $p$-value against a
bound is the level $\alpha$ at which the bound of the two-sided \IntervalLevel{} interval equals
it.
\item \textbf{Exact sign test.} $X$ counts the sessions whose value lies strictly on the
hypothesis side of the bound; a tie, or a session without a value, counts against. At the null's
boundary \SignNull, and $p = P(X \ge x)$ is computed exactly.
\item \textbf{Holm's step-down procedure} per family at family-wise $\alpha = \FwAlpha$,
one-sided, run once for each computation. A cell that cannot be tested enters with \PIsOne{}.
\end{itemize}
\ForAlex{Holm (1979) is not cited: no DOI is registered for it, and JSTOR refused a scripted
request, so its metadata could not be verified. Add the reference from your own copy.}

\runin{Equivalence in the cost family.} The author fixed the margin before any code existed. The
widest difference still called ``no measurable cost'' is a ratio between \MarginLow{} and
\MarginHigh, for C1, C2 and C3, on both hosts. Each cell has two one-sided nulls: the median ratio
is at most \MarginLow, and it is at least \MarginHigh. Each is tested at $\alpha$ by both
computations, and a computation's $p$-value for the cell is the larger of its two (two one-sided
tests~\cite{schuirmann1987tost}). Before Holm's adjustment, the BCa form passes exactly when the
two-sided \CiPct\% interval lies inside the margin. Each cost cell shows that interval and the
interval at its Holm level, \HolmLevel, where $j$ is its rank. A noisy cell, or one with too few
valid sessions, fails: a difference that is merely not significant never counts as equivalence.
In the cost family the sign test is the binding computation. At Holm's first threshold it allows
fewer sessions beyond a bound than the BCa needs: \SignAllowedLow{} beyond each bound at $R =
\CostCandMin$ and \SignAllowedHigh{} at $R = \CostCandMax$. With the acceleration at zero, as for
a median at an even $R$, a cost cell's BCa can pass $\alpha/\CostCellsMax$ only if
\ZzeroBound.

\runin{Superiority.} Each cell of B3 and of the mechanism family has one null, on the far side of
its bound from the hypothesis. For B3 the null is that the median $Q$ is at most
\FootprintMargin. For M1 and M3 it is that the median ratio is at most \MechBound, and for M2 that
it is at least \MechBound. The sign test counts
the sessions strictly on the hypothesis side. The BCa $p$-value is the level at which the lower
end of the interval equals the bound (the upper end for M2).

\runin{Replicates.} An exact paired test must be able to pass the family's first Holm threshold,
\DnineRule. That needs at least \CostDnineMin{} sessions in the cost family with
\CostCellsMax{} cells, \FootprintDnineMin{} for B3 and \MechDnineMin{} for the mechanism family.
B3 and the mechanism
family run $R_B = R_M = \RepsBM$, the reference count of the earlier
study~\cite{tsvetanov2026routing}, a design choice between that minimum and \RepsCap. No pilot can
size a superiority family, since nothing in the design sets its effect size. At that $R$, one
session on the wrong side of the bound still clears every Holm threshold of both families. Two
clear $\alpha/k$ only for $k \le \SupTwoK$, and three only for $k \le \SupThreeK$. There is no
sequential stopping.

\subsubsection{Sizing \texorpdfstring{\RC}{R\_C} from an A/A Pilot}\label{sec:pilot}

The pilot sizes $R$ for the cost family; it sets no margin. Every cost cell runs \PilotSessions{}
sessions with both arms in dedicated mode. The arms are two starts of the same binary and flags,
the second on a fixed port offset. So every window transition changes port, as in the A/B design.
The pilot holds no one-port data. For each cell, with \LogRatio{} the log of session $i$'s ratio:
\begin{enumerate}
\item the values are centred, $e_i = y_i - \mathrm{median}(y)$, so power is computed at a true
ratio of one;
\item $s$ is the standard deviation of the $e_i$. The bandwidth is Silverman's rule of
thumb~\cite{silverman2018density} divided by $s$: \Bandwidth. The spread used is \UpperSd, the
one-sided \UpperBoundPct\% upper confidence bound under normality. Here $q$ is the \ChiQuantile{}
quantile of the chi-square distribution with \PminusOne{} degrees of freedom, so $q = \ChiQ$ and
$s_U = \SUFactor\, s$;
\item one simulated run draws $R$ standardised values $e_i/s$ with replacement, adds to each an
independent normal draw with standard deviation $b$, multiplies by \SimScale{} and takes
exponentials;
\item the confirmatory code itself analyses each run, both computations at the margin, and the run
passes if both $p$-values are at most $\alpha/\CostCellsMax$;
\item the power at $R$ is the share of \NSim{} runs that pass (its Monte Carlo standard error at
\PowerTarget{} is \NSimSe).
\end{enumerate}
The candidates for $R$ are \CostCandidates: the values at which the sign test's allowed count
steps.
Runs share random numbers across candidates, and every stream is seeded from a frozen seed, so no
result depends on the order of computation. A cell is resolved if it has \PilotSessions{} valid
pilot sessions and power of at least \PowerTarget{} at $R = \CostCandMax$. \RC{} is the smallest
candidate at which every resolved cell reaches \PowerTarget. The pilot entry records \RC{} =
\CostR{} and \mC{} = \CostM. A cell not resolved leaves the confirmatory family, still runs at \RC, and is
reported as not resolved at $R \le \RunsMax$. For each protocol, backend and host whose C1, C2 and
C3 cells are all resolved, the product of their three powers is reported, not targeted.

\runin{Orders and seeds.} All sessions of one family's cells on one host run in one shuffled order,
so that drift spreads over the cells. The \FrozenSeeds{} seeds of the orders, of each family's
resamples and of the pilot were drawn once on L, in one revision-log entry before the code
freeze. Every integer that the lab journal or the engineering notes named was excluded, so no
seed of a development run could recur.

\subsection{Order of Freezes and Runs}\label{sec:order}

Each step starts only after the one before it is committed.
\begin{enumerate}
\item The pre-specification is frozen before any code of the server exists.
\item Engineering. Before it ends, one revision-log entry fixes every seed, and the analysis code
that sizes \RC{} is committed with its tests. That commit is \hex{\AnalysisCommit}; it pins numpy
\NumpyVersion{} on Python \PythonVersion.
\item The code freeze: one commit of the server, the generators, the holder, the harnesses, the
pins and the tests, with green sanitizer records and a passing test suite.
\item On the frozen binary, in dedicated mode only: the A/A pilot on each host, with its timer and
split parts.
\item The simulation that sizes \RC, on the pilot's data alone.
\item The pilot entry: one revision-log entry made from the pilot's and the simulation's outputs
alone, committed before any one-port window of the frozen code. Nothing in it is chosen.
\item After it: rule~E's choices, M2's rates, the footprint feasibility windows, \Kbase{}, and
every confirmatory, secondary and hard-case run, each family in its own order.
\end{enumerate}

\runin{The pilot's parts.} The timer part runs \TimerRuns{} single connections per backend whose
PROXY header stays incomplete, against the dedicated HTTP/1.1 port, and records how late the loop
handles \Thdr. Per host, $G$ is the smallest whole number of milliseconds above the largest
lateness. It sets the guard band of the hard case whose first byte arrives just after or just
before \Tfb. If a host has no valid run, or $G$ is not below \Tfb, that case is not run there, and
the paper says why. The split part sends an HTTP/1.1 request split after its first byte, with gaps
of \SplitGaps~ms, \SplitReps{} times each. The smallest gap at which every replicate on every
backend shows the split read becomes the gap of the split hard cases.
\TODO{analysis/macros.py emits no macro for $G_L$, $G_W$, the split gap, $\lambda$ per C3 cell,
M2's rates or rule~E's choices. Add them before the results are written.}

\runin{Rule~E.} Three choices per backend are fixed by a rule: the default detection mode, the
receive form on IOCP, and the relay copy. Both options of each are built into the frozen binary as
flag values, because a change of the binary's layout alone can move timings. An option replaces
the proposed default (replay, the zero-byte receive, user-space buffers) only under two conditions.
Every session of at least \RuleESessions{} development sessions of HTTP/1.1 churn must favour it,
and its operations per connection must not be higher. That count of sessions is the smallest at
which an exact sign test can reach $p = \RuleEP$, below $\alpha$. The receive form that posts a
buffer at once breaks B2's memory bound, so rule~E cannot choose it (a decision the revision log
records), and the zero-byte form stands. The other two choices exist only in one-port mode and are
made after the pilot entry, before any confirmatory window.

\runin{Development data.} Engineering ran development windows. They are never results; they are
cited only in the revision log and the lab journal, and disclosed here. Before the code freeze they
could time dedicated mode against anything, and one-port mode against a competitor or against
another one-port option. No window, on any code, timed one-port mode against dedicated mode before
the pilot entry. A dedicated window and a one-port window against the same competitor, at the same
protocol, backend and host, would still give an indirect ratio of the two modes. The revision log
lists every such pair when it runs; at the time of writing it lists none, and no rule uses one.
\ForAlex{Confirm that the revision log still lists no such pair before submission.}

\runin{A change once the pilot has started.} Before the simulation starts, a change forced by a
named failure needs new records and a new pilot, at most once. The named failures are a failing
test, a sanitizer record that is not green, the gate refusing a build, and a crash of a
first-party program. Any other change keeps the pilot and gets new records. Every confirmatory
window of the old build is archived, the runs start again on the new build, and the change is
recorded as a deviation. No one-port window of either build
runs before the pilot entry.

\subsection{Sanitizer Gate and Declared Gaps}\label{sec:gate}

A number is citable only if every measured build compiled first-party inputs that a green
sanitizer record covers. On L the records use clang \ClangVersion{} with AddressSanitizer and
UndefinedBehaviorSanitizer together, ThreadSanitizer, and MemorySanitizer with an instrumented C++
standard library. On W they use MSVC \MsvcVersion{} with AddressSanitizer, one compiler because the
gate matches the compiler. Each record covers the server, \opgen, \opcase, \ophold{} and the test
suite by a hash of their inputs. The suite runs the detection table at every split and the whole
table of hard cases. It runs every mode of the binary, both detection modes with their peek paths,
and both forms of the check of rule (b). It also runs both dispatch modes, TLS handshakes, h2
sessions, the MQTT, SSH and SMTP exchanges, the counters and the holder. Under ThreadSanitizer, two
workers share one listener. On
L both Linux backends run, on W, IOCP. The gate refuses a measured build unless green records on
the same host cover every first-party target, with the same compiler, configuration, pins and
fetched archives. The Go and Rust harnesses get AddressSanitizer and ThreadSanitizer records.

Declared gaps are listed in the repository's coverage file:
\begin{itemize}
\item OpenSSL's assembly, which no sanitizer instruments; its MemorySanitizer build turns assembly
off, so that sanitizer covers the C fallbacks only;
\item the function check of UndefinedBehaviorSanitizer, left out of OpenSSL alone, which calls
typed functions through generic pointer types by design;
\item the kernel's writes into io\_uring's rings and buffers, which MemorySanitizer cannot see;
every structure a raw call fills is initialised, and exactly the bytes a completion reports are
unpoisoned;
\item the competitors' binaries and runtimes, which are pinned, not instrumented;
\item the harnesses: no MemorySanitizer for the Go and Rust harnesses, no sanitizer for the Java
harnesses, and no UndefinedBehaviorSanitizer form for Go or Rust;
\item on Windows, ThreadSanitizer and MemorySanitizer, which the rule does not require there.
\end{itemize}
OpenSSL itself ran under ThreadSanitizer: engineering built it instrumented and the suite passed,
so that gap was withdrawn, as the revision log records.

\subsection{Secondary Analyses and Cells Not Run}\label{sec:secondary}

Secondary analyses decide nothing and are in no Holm family. Cells that need sessions run at $R =
\ReplicatesSecondary$, and each result is reported with its \CiPct\% interval where it has one.
They are these.
\begin{itemize}
\item The CPU per exchange and resident memory of every cost cell, and the tail of the time to
first byte.
\item The cost cells that the pilot did not resolve.
\item SSH in the cost family (\SshSecCells{} cells). The server's identification line waits for
the client's in one-port mode but is sent at accept in dedicated mode, so the message order
differs by mode.
\item A mixed-protocol cell per backend (\MixedSecCells{} cells): HTTP/1.1 churn beside a fixed
background of \NBg{} kept-alive TLS, \NBg{} kept-alive MQTT and \NBg{} silent connections.
\item TLS with session resumption and TLS with ALPN \texttt{h2} (\TlsSecCells{} cells).
\item Two server cores (\TwoCoreSecCells{} cells), and a \texttt{SO\_REUSEPORT} group against the
shared listener (\ReuseportSecCells{} cells).
\item The server's relay on io\_uring against the proxies (\RelaySecCells{} cells).
\item Two accept and receive forms on IOCP (\IocpSecCells{} cells). One is the posted form, which
breaks B2's memory bound; the paper states that beside it.
\item The time to first byte at a fixed load for M1 and M3 (\TtfbSecCells{} cells).
\item B3 with the server in its other detection mode (\OtherModeSecCells{} cells).
\item Operations and copies per connection.
\item The competitors' outcomes on the hard cases they cover, at matched timers and at their
defaults, \CompReps{} replicates each. A system with no timer is observed for at most
\TObsS{}~s.
\item The distributions of B2.
\end{itemize}

On L, one untimed window per cost cell checks the server's system-call counters against the
kernel's count of the same calls' entry tracepoints, read with \texttt{perf stat}. \texttt{perf
trace -s} stays beside it as a cross-check, and its shortfall is reported beside the counts.

Two kinds of cell are reported as not run, each with its reason. TLS with session resumption
(\ResumptionCells{} cells) is not run. The frozen settings of every endpoint exclude session tickets
and a session cache, and TLS~1.3 resumes only from a ticket. An M2 cell whose rate its
rule cannot compute, because an arm has no valid development session, has no rate. It runs no
window, enters Holm with \PIsOne{}, and keeps the family's size; it is reported as a difference not
shown.

\subsection{What Counts against the Claims}\label{sec:against}

\runin{Cost.} A cell that fails is reported with its interval as ``equivalence not shown''. It is
called a measured cost only where its \CiPct\% interval lies wholly outside the margin. ``No
measurable cost'' is claimed only as equivalence within the margin, and for load of one protocol.
It is claimed only for a protocol, backend and host where C1, C2 and C3 all pass. The joint power
and the intervals at Holm's level stand beside it.

\runin{Robustness.} B1 and B2 are not narrowed: a failure in any run of the frozen code is reported
as a failure, and the robustness claim is not made. B3 is narrowed to the competitors and cases
where it passes. A pass is written as ``the competitor held at least \FootprintMargin{} times the
server's counted footprint per pending connection''. It always comes with the difference in bytes
and its interval, and with WL7's exclusions and residuals named. ``Loss'' is used only where the \CiPct\%
interval of $Q$ lies wholly below one.

\runin{Mechanism.} M1 and M2 are reported per backend as measured; a cell that fails shows only
that a difference was not shown. M3 is narrowed to the proxies and protocols where it passes.

\runin{Narrowing order.} The cost family leads the paper. If a family cannot be won cleanly, the
claims narrow in a fixed order. M3 narrows first, to the proxies and protocols the server beats.
B3 follows, to the competitors and cases it beats. The cost family narrows last, per backend.

\runin{Analysis.} The analysis computes every decision from the archived runs and writes the
paper's numbers into a generated macro file. Raw outputs are archived with their SHA-256. Nothing
is run again to change a decision.
